\documentclass[12pt, twoside]{article}
%\documentclass[12pt, twoside]{article}
\usepackage[letterpaper, margin=1in, headsep=0.2in]{geometry}
\setlength{\headheight}{0.6in}
%\usepackage[english]{babel}
\usepackage[utf8]{inputenc}
\usepackage{microtype}
\usepackage{amsmath}
\usepackage{amssymb}
%\usepackage{amsfonts}
\usepackage[nomessages]{fp} %\FPeval{\var-name}{2*sin(pi/6)}
\usepackage{siunitx} %units in math. eg 20\milli\meter
\usepackage{yhmath} % for arcs, overparenth command
\usepackage{tikz} %graphics
\usetikzlibrary{quotes, angles, arrows, arrows.meta}
\usepackage{pgfplots}
\usepgfplotslibrary{fillbetween}
\pgfplotsset{compat=1.16}
\usepackage{graphicx} %consider setting \graphicspath{{images/}}
\usepackage{parskip} %no paragraph indent
\usepackage{enumitem}
\usepackage{multicol}
\usepackage{venndiagram}

\usepackage{fancyhdr}
\pagestyle{fancy}
\fancyhf{}
\renewcommand{\headrulewidth}{0pt} % disable the underline of the header
\raggedbottom
\hfuzz=2mm %suppresses overfull box warnings

\usepackage{hyperref}

\fancyhead[LO]{La Scuola d'Italia / Dr.\ Huson / IB Math 11 \\* 7 October 2026}
\fancyhead[RO]{\fbox{\begin{minipage}{6cm}\footnotesize Nickname:\\[0.5cm]\end{minipage}}}
\fancyhead[LE]{\fbox{\begin{minipage}{6cm}\footnotesize Nickname:\\[0.5cm]\end{minipage}}}
\fancyfoot[C]{\thepage}

\begin{document}
\subsubsection*{1.10 Early Finishers: Sequences and Compound Interest}

\emph{For students who finish the quiz early. A calculator is allowed.
Show your working: write the formula you use, then substitute.}

\begin{enumerate}[itemsep=0.15cm]

%--- Problem 1: Arithmetic sequence from two conditions; count of terms below 500; arithmetic terms forming a geometric sequence ---
%--- Source: Original (freshly authored, AA SL 1.2, 1.3); extension of 1-10Q_Sequences P1-P3 ---
\item[1.]

In an arithmetic sequence, $u_{3} + u_{8} = 48$ and $u_{12} = 5\,u_{2}$.

\begin{enumerate}[itemsep=0.1cm]
  \item Find the first term $u_{1}$ and the common difference $d$.
  \vfill
  \item Find the number of terms of this sequence that are less than
  $500$.
  \vfill
\end{enumerate}

A different arithmetic sequence has first term $3$ and common difference
$d$, where $d \neq 0$. Its $1^{\text{st}}$, $2^{\text{nd}}$ and
$4^{\text{th}}$ terms are the first three terms of a geometric sequence.

\begin{enumerate}[resume, itemsep=0.1cm]
  \item Find the value of $d$ and the common ratio of the geometric
  sequence.
\end{enumerate}

\vfill

\newpage

%--- Problem 2: Compound interest -- doubling time; unknown rate; first year one investment overtakes another ---
%--- Source: Original (freshly authored, AA SL 1.4); extension of 1-7CW_Financial, 1-8CW_Inflation ---
\item[2.]

Sofia invests \$$5000$ in an account that pays a nominal annual interest
rate of $4.8\%$, compounded monthly.

\begin{enumerate}[itemsep=0.1cm]
  \item Find the least number of whole years after which Sofia's
  investment is worth more than \$$10\,000$.
  \vfill
\end{enumerate}

At the same time, Luca invests \$$6000$ in an account that pays a
nominal annual interest rate of $r\%$, compounded quarterly. After
$10$ years Luca's investment is worth \$$9000$.

\begin{enumerate}[resume, itemsep=0.1cm]
  \item Find the value of $r$.
  \vfill
  \item Find the first whole year at the end of which Sofia's
  investment is worth more than Luca's.
\end{enumerate}

\vfill

\end{enumerate}
\end{document}
